\chapter*{Abstract}
\addcontentsline{toc}{chapter}{Abstract}

SOC teams receive heterogeneous alerts from host, network, runtime, and metric sources, yet most tools leave correlation, investigation, and remediation as manual, fragmented tasks. This final-year project presents \project{} (\projectlong{}), an open SOAR/SOC automation platform that closes this gap by turning raw Elasticsearch telemetry into normalized, enriched, correlated incidents and supporting AI-assisted, human-approved Ansible remediation with post-fix verification.

ARIA ingests Wazuh, Suricata/Filebeat, Falco, Telegraf, and generic alerts; deduplicates and filters noise; enriches events with GeoIP and MITRE ATT\&CK context; groups related alerts into incidents; and uses a FastAPI backend plus Next.js dashboard to expose triage, investigation, remediation, and monitoring workflows. An AI response layer generates summaries, risk scores, narratives, and playbooks, while approval gates, admin-secret protection, safety checks, and audit logging keep remediation controlled.

Validation confirms a working multi-server deployment, 910 passing automated tests, backend/frontend build success, and honest limits around single-node SQLite storage and internal-trusted authentication.

\textbf{Keywords:} SOAR, SOC automation, FastAPI, Next.js, Elasticsearch, Redis, SQLite, Ansible, MITRE ATT\&CK, Sigma, AI-assisted incident response.

\vspace{1cm}
\chapter*{Summary}
\addcontentsline{toc}{chapter}{Summary}

This final-year project focuses on the design and implementation of \project{}, also called \projectlong{}, a platform for automating security and system operations. The project aims to reduce fragmentation in SOC and system-team workflows by connecting alert ingestion, enrichment, incident correlation, AI-assisted investigation, human approval, Ansible remediation, post-fix verification, archiving, monitoring, and visualization into one coherent solution.

The implemented platform collects data from Elasticsearch, including Wazuh, Suricata/Filebeat, Falco, Telegraf, and generic alerts. It normalizes security and observability signals, removes duplicates, filters noise, adds GeoIP and MITRE ATT\&CK context, correlates events, and stores operational data in SQLite while using Redis for cursors, cache, deduplication, and metrics. The FastAPI backend exposes features through REST and WebSocket, while the Next.js dashboard provides the operational interfaces.

\textbf{Keywords:} SOAR, SOC, FastAPI, Next.js, Elasticsearch, Redis, SQLite, Ansible, MITRE ATT\&CK, Sigma, incident response, system monitoring, observability.

